\subsection{Homology: the two complexes-of-complexes conventions}\label{homology-the-two-complexes-of-complexes-conventions}

This editorial note continues PROOFS\_2799\_4336.md. Primary source: Stacks Project, homology.tex, official commit a04446e57ec1fbc252a871afcec7752fb2807b14, especially remark-double-complex-complex-of-complexes-first (4304--4475) and remark-double-complex-complex-of-complexes-second (4477--4651). The pinned cumulative source is commit b4b30d521cfbe300860903ff788059627da6a434. No novelty is claimed.

\subsubsection{1. Twelve reports represent four already corrected findings}\label{twelve-reports-represent-four-already-corrected-findings}

{\def\LTcaptype{none} % do not increment counter
\begin{longtable}[]{@{}
  >{\raggedright\arraybackslash}p{(\linewidth - 4\tabcolsep) * \real{0.3333}}
  >{\raggedright\arraybackslash}p{(\linewidth - 4\tabcolsep) * \real{0.3333}}
  >{\raggedright\arraybackslash}p{(\linewidth - 4\tabcolsep) * \real{0.3333}}@{}}
\toprule\noalign{}
\begin{minipage}[b]{\linewidth}\raggedright
Group
\end{minipage} & \begin{minipage}[b]{\linewidth}\raggedright
Source locus
\end{minipage} & \begin{minipage}[b]{\linewidth}\raggedright
Existing correction
\end{minipage} \\
\midrule\noalign{}
\endhead
\bottomrule\noalign{}
\endlastfoot
HOMOLOGY-RECON-026 & 4435 & MC-STK-ERR-0774: the component s\^{}\{p,q\} lands in B\^{}\{p,q\}. \\
HOMOLOGY-RECON-027 & 4518 & MC-STK-ERR-0775: h\^{}\{p,q\} is the degree-q component of h\^{}p. \\
HOMOLOGY-RECON-028 & 4581 & MC-STK-ERR-0776: the first-index connecting map lands in A{[}1,0{]}. \\
HOMOLOGY-RECON-029 & 4608--4609 & MC-STK-ERR-0777: the component lands in B\^{}\{p,q\}, and the families are s\^{}p and pi\^{}p. \\
\end{longtable}
}

Each group has one original p02 report, one later p02 report explicitly withdrawn by its producer as a duplicate, and one French report. The withdrawal is preserved. The twelve physical reports do not become twelve corrections: there are four existing units, containing five existing source operations. All current postimages are checked exactly in REVIEW\_PART04.json.

\subsubsection{2. Homotopies and shifts in both conventions}\label{homotopies-and-shifts-in-both-conventions}

Work in the source's additive category with countable direct sums. The double differentials commute and the total differential on A\^{}\{p,q\} is d\_1\textsuperscript{\{p,q\}+(-1)}p d\_2\^{}\{p,q\}.

In the first convention the outer degree is q. A homotopy h\^{}q between f and g has components \[
h^{p,q}:A^{p,q}\longrightarrow B^{p,q-1}.
\] The outer homotopy equation and the assertion that h\^{}q is an inner cochain map are respectively \[
f^{p,q}-g^{p,q}
=d_{B,2}^{p,q-1}h^{p,q}+h^{p,q+1}d_{A,2}^{p,q},
\] \[
d_{B,1}^{p,q-1}h^{p,q}=h^{p+1,q}d_{A,1}^{p,q}.
\] Define H on the summand A\^{}\{p,q\} of Tot\^{}n(A), n=p+q, to be (-1)\^{}p h\^{}\{p,q\}, with target the indicated summand of Tot\^{}\{n-1\}(B). The full component of D\_B H+H D\_A is \[
\begin{aligned}
&(-1)^p d_{B,1}^{p,q-1}h^{p,q}
 +(-1)^{2p}d_{B,2}^{p,q-1}h^{p,q}\\
&\hspace{1em}
 +(-1)^{p+1}h^{p+1,q}d_{A,1}^{p,q}
 +(-1)^{2p}h^{p,q+1}d_{A,2}^{p,q}.
\end{aligned}
\] The two terms with d\_1 have opposite coefficients and equal composites. The other two terms give f\textsuperscript{\{p,q\}-g}\{p,q\}. Thus Tot sends this homotopy to a homotopy with the asserted sign.

The outer shift by k is {[}0,k{]}. By Section 5 of the preceding note, the comparison \[
\gamma_{0,k}:\operatorname{Tot}(B)[k]
\longrightarrow\operatorname{Tot}(B[0,k])
\] on the original summand B\^{}\{p,q\} is (-1)\^{}\{pk\}. If f=g, the homotopy equation says exactly that h:A -\textgreater{} B{[}0,-1{]} is a cochain map in the outer direction: -d\_\{B,2\}h=h d\_\{A,2\}. It is also a map in the inner direction by the displayed inner equation. Composing Tot(h) with gamma\_\{0,-1\}\^{}\{-1\} multiplies its component by (-1)\textsuperscript{\{-p\}=(-1)}p and therefore gives precisely H as a map Tot(A) -\textgreater{} Tot(B){[}-1{]}.

In the second convention the outer degree is p.~The homotopy h\^{}p has degree-q component \[
h^{p,q}:A^{p,q}\longrightarrow B^{p-1,q}.
\] Consequently it is not a component of h\^{}q. Its equations are \[
f^{p,q}-g^{p,q}
=d_{B,1}^{p-1,q}h^{p,q}+h^{p+1,q}d_{A,1}^{p,q},
\qquad
d_{B,2}^{p-1,q}h^{p,q}=h^{p,q+1}d_{A,2}^{p,q}.
\] Put H\textbar{}\emph{\{A\textsuperscript{\{p,q\}\}=h}\{p,q\}, with no additional sign. Then \[
\begin{aligned}
(D_BH+HD_A)|_{A^{p,q}}
={}&d_{B,1}^{p-1,q}h^{p,q}
 +(-1)^{p-1}d_{B,2}^{p-1,q}h^{p,q}\\
 &+h^{p+1,q}d_{A,1}^{p,q}
 +(-1)^p h^{p,q+1}d_{A,2}^{p,q}.
\end{aligned}
\] The d\_2 terms cancel by the inner-map equation, and the d\_1 terms give f-g. The shift comparison gamma}\{k,0\} has component (-1)\^{}\{p\cdot0\}=1. For f=g, the outer homotopy is a map to B{[}-1,0{]}, and Tot(h) followed by gamma\_\{-1,0\}\^{}\{-1\} is exactly H. These calculations verify the component and degree corrections above and preserve the different signs in the two conventions.

\subsubsection{3. Split boundaries and the two remaining projection-index errors}\label{split-boundaries-and-the-two-remaining-projection-index-errors}

Let \[
0\longrightarrow A^{\bullet,\bullet}
\xrightarrow{u}B^{\bullet,\bullet}
\xrightarrow{v}C^{\bullet,\bullet}\longrightarrow0
\] have the stated splittings. In each bidegree they give \[
v s=1,\quad \pi u=1,\quad \pi s=0,\quad u\pi+s v=1.
\] The induced total splitting is the direct sum of these component maps, so these same identities hold on each total term.

First suppose the splittings are morphisms of the p-complexes, indexed by q: \[
s^q:C^{\bullet,q}\longrightarrow B^{\bullet,q},\qquad
\pi^q:B^{\bullet,q}\longrightarrow A^{\bullet,q}.
\] Their degree-p components have types s\textsuperscript{\{p,q\}:C}\{p,q\} -\textgreater{} B\^{}\{p,q\} and pi\textsuperscript{\{p,q\}:B}\{p,q\} -\textgreater{} A\^{}\{p,q\}, confirming MC-STK-ERR-0774. The boundary in the outer complex is \[
\boxed{\quad
\delta^q=\pi^{q+1}\,d_{B,2}^{\,q}\,s^q:
C^{\bullet,q}\longrightarrow A^{\bullet,q+1}.
\quad}                                                     \tag{3.1}
\] The middle map has codomain B\^{}\{bullet,q+1\}, which is the domain of pi\^{}\{q+1\}, not pi\^{}q. The source's line 4468 incorrectly prints pi\^{}q. Equation (3.1) is also forced by the source's correctly indexed component formula at line 4472.

On C\^{}\{p,q\}, the boundary of the total split sequence is \[
\begin{aligned}
\Delta^{p,q}
={}&\pi^{p+1,q}d_{B,1}^{p,q}s^{p,q}
 +(-1)^p\pi^{p,q+1}d_{B,2}^{p,q}s^{p,q}\\
={}&0+(-1)^p\pi^{p,q+1}d_{B,2}^{p,q}s^{p,q}.
\end{aligned}                                               \tag{3.2}
\] Indeed d\_\{B,1\}\textsuperscript{\{p,q\}s}\{p,q\}=s\textsuperscript{\{p+1,q\}d\_\{C,1\}}\{p,q\}, because s\^{}q is an inner cochain map, and pi\textsuperscript{\{p+1,q\}s}\{p+1,q\}=0. The remaining term is exactly the component of Tot(delta) followed by gamma\_\{0,1\}\^{}\{-1\}, whose component on A\^{}\{p,q+1\} is (-1)\^{}p. Thus the source's totalization comparison holds, with all domains, codomains and signs displayed.

Next suppose the splittings are morphisms of the q-complexes, indexed by p: \[
s^p:C^{p,\bullet}\longrightarrow B^{p,\bullet},\qquad
\pi^p:B^{p,\bullet}\longrightarrow A^{p,\bullet}.
\] The components are degree-q components of these families. They have the same bidegree types as above. The boundary now is \[
\boxed{\quad
\delta^p=\pi^{p+1}\,d_{B,1}^{\,p}\,s^p:
C^{p,\bullet}\longrightarrow A^{p+1,\bullet}.
\quad}                                                     \tag{3.3}
\] The source's line 4641 incorrectly prints pi\^{}p.~The next-degree projection pi\^{}\{p+1\} is required by the codomain of d\_\{B,1\} and by the correctly indexed component formula at line 4645. The map delta in this convention has target A{[}1,0{]}, as the existing correction MC-STK-ERR-0776 states. It is a morphism of double complexes: it commutes with the inner q-differential because each factor in (3.3) is a map of q-complexes; in the p-direction its cochain equation is the split-boundary identity of Section 3 in the preceding note, with the negative differential on the shifted target.

For the total boundary, the first line of (3.2) still applies, but now the second term vanishes: \[
\pi^{p,q+1}d_{B,2}^{p,q}s^{p,q}
=\pi^{p,q+1}s^{p,q+1}d_{C,2}^{p,q}=0.
\] What remains is \[
\Delta^{p,q}=\pi^{p+1,q}d_{B,1}^{p,q}s^{p,q},
\] the component of (3.3). Since gamma\_\{1,0\}\^{}\{-1\} has all components equal to identities, this is precisely the other totalization comparison. The corrections of the two projection indices restore the displayed composite types; they do not change the correctly printed component calculations or the totalization conclusions.

The two new operations form HOMOLOGY-FIND-004. No received report in the bound Homology inventory has a parsed interval containing either 4468 or 4641. The inventory and earlier-operation checks are retained in CLAIM\_PROPAGATION\_PART04.json.

\subsubsection{4. Receiving use and limits}\label{receiving-use-and-limits}

The exact top-level TeX reference search for these two remark labels at the pinned cumulative commit found their uses in derived.tex, remark-homotopy-double, lines 3858 and 3876. The surrounding passage 3828--3900 was read. It uses the two totalization functors on the two homotopy categories. Sections 2 and 3 above verify the relevant homotopy, shift and split-boundary comparisons with the original orientations intact. No additional body correction in that receiving passage is required. The foundational construction of its triangulated categories is not claimed to have been reaudited in full here.

Primary-source reading continued through line 4680, where the filtered-object definition is still in progress. Reports through the end of the second double-complex remark, line 4651, have been adjudicated. The report at 4675--4677 has been read but is left for the next coherent filtration review. Next source line: 4681.

All new edits remain exact local previews, with full inverse checks. The accepted source and all inherited additions remain unchanged.
